% Options for packages loaded elsewhere
\PassOptionsToPackage{unicode}{hyperref}
\PassOptionsToPackage{hyphens}{url}
%
\documentclass[
]{article}
\usepackage{lmodern}
\usepackage{amssymb,amsmath}
\usepackage{ifxetex,ifluatex}
\ifnum 0\ifxetex 1\fi\ifluatex 1\fi=0 % if pdftex
  \usepackage[T1]{fontenc}
  \usepackage[utf8]{inputenc}
  \usepackage{textcomp} % provide euro and other symbols
\else % if luatex or xetex
  \usepackage{unicode-math}
  \defaultfontfeatures{Scale=MatchLowercase}
  \defaultfontfeatures[\rmfamily]{Ligatures=TeX,Scale=1}
\fi
% Use upquote if available, for straight quotes in verbatim environments
\IfFileExists{upquote.sty}{\usepackage{upquote}}{}
\IfFileExists{microtype.sty}{% use microtype if available
  \usepackage[]{microtype}
  \UseMicrotypeSet[protrusion]{basicmath} % disable protrusion for tt fonts
}{}
\makeatletter
\@ifundefined{KOMAClassName}{% if non-KOMA class
  \IfFileExists{parskip.sty}{%
    \usepackage{parskip}
  }{% else
    \setlength{\parindent}{0pt}
    \setlength{\parskip}{6pt plus 2pt minus 1pt}}
}{% if KOMA class
  \KOMAoptions{parskip=half}}
\makeatother
\usepackage{xcolor}
\IfFileExists{xurl.sty}{\usepackage{xurl}}{} % add URL line breaks if available
\IfFileExists{bookmark.sty}{\usepackage{bookmark}}{\usepackage{hyperref}}
\hypersetup{
  pdftitle={Metabodecon Models},
  hidelinks,
  pdfcreator={LaTeX via pandoc}}
\urlstyle{same} % disable monospaced font for URLs
\usepackage[margin=1in]{geometry}
\usepackage{color}
\usepackage{fancyvrb}
\newcommand{\VerbBar}{|}
\newcommand{\VERB}{\Verb[commandchars=\\\{\}]}
\DefineVerbatimEnvironment{Highlighting}{Verbatim}{commandchars=\\\{\}}
% Add ',fontsize=\small' for more characters per line
\usepackage{framed}
\definecolor{shadecolor}{RGB}{248,248,248}
\newenvironment{Shaded}{\begin{snugshade}}{\end{snugshade}}
\newcommand{\AlertTok}[1]{\textcolor[rgb]{0.94,0.16,0.16}{#1}}
\newcommand{\AnnotationTok}[1]{\textcolor[rgb]{0.56,0.35,0.01}{\textbf{\textit{#1}}}}
\newcommand{\AttributeTok}[1]{\textcolor[rgb]{0.77,0.63,0.00}{#1}}
\newcommand{\BaseNTok}[1]{\textcolor[rgb]{0.00,0.00,0.81}{#1}}
\newcommand{\BuiltInTok}[1]{#1}
\newcommand{\CharTok}[1]{\textcolor[rgb]{0.31,0.60,0.02}{#1}}
\newcommand{\CommentTok}[1]{\textcolor[rgb]{0.56,0.35,0.01}{\textit{#1}}}
\newcommand{\CommentVarTok}[1]{\textcolor[rgb]{0.56,0.35,0.01}{\textbf{\textit{#1}}}}
\newcommand{\ConstantTok}[1]{\textcolor[rgb]{0.00,0.00,0.00}{#1}}
\newcommand{\ControlFlowTok}[1]{\textcolor[rgb]{0.13,0.29,0.53}{\textbf{#1}}}
\newcommand{\DataTypeTok}[1]{\textcolor[rgb]{0.13,0.29,0.53}{#1}}
\newcommand{\DecValTok}[1]{\textcolor[rgb]{0.00,0.00,0.81}{#1}}
\newcommand{\DocumentationTok}[1]{\textcolor[rgb]{0.56,0.35,0.01}{\textbf{\textit{#1}}}}
\newcommand{\ErrorTok}[1]{\textcolor[rgb]{0.64,0.00,0.00}{\textbf{#1}}}
\newcommand{\ExtensionTok}[1]{#1}
\newcommand{\FloatTok}[1]{\textcolor[rgb]{0.00,0.00,0.81}{#1}}
\newcommand{\FunctionTok}[1]{\textcolor[rgb]{0.00,0.00,0.00}{#1}}
\newcommand{\ImportTok}[1]{#1}
\newcommand{\InformationTok}[1]{\textcolor[rgb]{0.56,0.35,0.01}{\textbf{\textit{#1}}}}
\newcommand{\KeywordTok}[1]{\textcolor[rgb]{0.13,0.29,0.53}{\textbf{#1}}}
\newcommand{\NormalTok}[1]{#1}
\newcommand{\OperatorTok}[1]{\textcolor[rgb]{0.81,0.36,0.00}{\textbf{#1}}}
\newcommand{\OtherTok}[1]{\textcolor[rgb]{0.56,0.35,0.01}{#1}}
\newcommand{\PreprocessorTok}[1]{\textcolor[rgb]{0.56,0.35,0.01}{\textit{#1}}}
\newcommand{\RegionMarkerTok}[1]{#1}
\newcommand{\SpecialCharTok}[1]{\textcolor[rgb]{0.00,0.00,0.00}{#1}}
\newcommand{\SpecialStringTok}[1]{\textcolor[rgb]{0.31,0.60,0.02}{#1}}
\newcommand{\StringTok}[1]{\textcolor[rgb]{0.31,0.60,0.02}{#1}}
\newcommand{\VariableTok}[1]{\textcolor[rgb]{0.00,0.00,0.00}{#1}}
\newcommand{\VerbatimStringTok}[1]{\textcolor[rgb]{0.31,0.60,0.02}{#1}}
\newcommand{\WarningTok}[1]{\textcolor[rgb]{0.56,0.35,0.01}{\textbf{\textit{#1}}}}
\usepackage{longtable,booktabs}
% Correct order of tables after \paragraph or \subparagraph
\usepackage{etoolbox}
\makeatletter
\patchcmd\longtable{\par}{\if@noskipsec\mbox{}\fi\par}{}{}
\makeatother
% Allow footnotes in longtable head/foot
\IfFileExists{footnotehyper.sty}{\usepackage{footnotehyper}}{\usepackage{footnote}}
\makesavenoteenv{longtable}
\usepackage{graphicx}
\makeatletter
\def\maxwidth{\ifdim\Gin@nat@width>\linewidth\linewidth\else\Gin@nat@width\fi}
\def\maxheight{\ifdim\Gin@nat@height>\textheight\textheight\else\Gin@nat@height\fi}
\makeatother
% Scale images if necessary, so that they will not overflow the page
% margins by default, and it is still possible to overwrite the defaults
% using explicit options in \includegraphics[width, height, ...]{}
\setkeys{Gin}{width=\maxwidth,height=\maxheight,keepaspectratio}
% Set default figure placement to htbp
\makeatletter
\def\fps@figure{htbp}
\makeatother
\setlength{\emergencystretch}{3em} % prevent overfull lines
\providecommand{\tightlist}{%
  \setlength{\itemsep}{0pt}\setlength{\parskip}{0pt}}
\setcounter{secnumdepth}{-\maxdimen} % remove section numbering
\usepackage{fontspec}
\usepackage{etoolbox}
\usepackage{fvextra}
\usepackage{needspace}
\definecolor{shadecolor}{RGB}{232,232,232}
\fvset{breaklines=true,breakanywhere=true}
\BeforeBeginEnvironment{Shaded}{\Needspace{10\baselineskip}}
\BeforeBeginEnvironment{Shaded}{\vspace{0.5em}}
\DefineVerbatimEnvironment{Highlighting}{Verbatim}{breaklines=true,breakanywhere=true,commandchars=\\\{\}}
\RecustomVerbatimEnvironment{verbatim}{Verbatim}{breaklines=true,breakanywhere=true}

\title{Metabodecon Models}
\author{}
\date{\vspace{-2.5em}}

\begin{document}
\maketitle

This article demonstrates how to use \texttt{fit\_mdm()} and
\texttt{benchmark()} to build, tune and evaluate classification models
on deconvoluted NMR spectra. We use the simulated \texttt{sim2} dataset
shipped with the package, split it into training and test sets, run a
small grid search to find good preprocessing parameters, and then
evaluate on held-out data.

\hypertarget{load-spectra}{%
\section{Load spectra}\label{load-spectra}}

\texttt{sim2} contains 100 simulated 1D NMR spectra split evenly into
groups \texttt{A} and \texttt{B}. Five out of 25 peaks per spectrum
differ between groups by 10\% in area. The group labels are attached as
an attribute. For details on how the dataset was constructed, see
\texttt{?sim2}.

\begin{Shaded}
\begin{Highlighting}[]
\KeywordTok{library}\NormalTok{(metabodecon)}
\KeywordTok{data}\NormalTok{(sim2)}
\NormalTok{spectra \textless{}{-}}\StringTok{ }\NormalTok{sim2}
\NormalTok{group \textless{}{-}}\StringTok{ }\KeywordTok{attr}\NormalTok{(sim2, }\StringTok{"group"}\NormalTok{)}
\NormalTok{n \textless{}{-}}\StringTok{ }\KeywordTok{length}\NormalTok{(spectra)}
\end{Highlighting}
\end{Shaded}

The following figure shows four spectra from each group overlaid on top
of each other. The subtle area differences are hard to spot visually.

\begin{Shaded}
\begin{Highlighting}[]
\NormalTok{idx\_A \textless{}{-}}\StringTok{ }\KeywordTok{which}\NormalTok{(group }\OperatorTok{==}\StringTok{ "A"}\NormalTok{)[}\DecValTok{1}\OperatorTok{:}\DecValTok{4}\NormalTok{]}
\NormalTok{idx\_B \textless{}{-}}\StringTok{ }\KeywordTok{which}\NormalTok{(group }\OperatorTok{==}\StringTok{ "B"}\NormalTok{)[}\DecValTok{1}\OperatorTok{:}\DecValTok{4}\NormalTok{]}
\KeywordTok{plot\_spectra}\NormalTok{(spectra[}\KeywordTok{c}\NormalTok{(idx\_A, idx\_B)])}
\end{Highlighting}
\end{Shaded}

\includegraphics[width=1\linewidth]{MDM_files/figure-latex/plot-simulated-1}

\hypertarget{traintest-split}{%
\section{Train/test split}\label{traintest-split}}

We use two thirds of the data for training and one third for testing.

\begin{Shaded}
\begin{Highlighting}[]
\KeywordTok{set.seed}\NormalTok{(}\DecValTok{1}\NormalTok{)}
\NormalTok{n\_train \textless{}{-}}\StringTok{ }\KeywordTok{round}\NormalTok{(}\DecValTok{2} \OperatorTok{/}\StringTok{ }\DecValTok{3} \OperatorTok{*}\StringTok{ }\NormalTok{n)}
\NormalTok{train\_idx \textless{}{-}}\StringTok{ }\KeywordTok{sort}\NormalTok{(}\KeywordTok{sample}\NormalTok{(n, n\_train))}
\NormalTok{test\_idx \textless{}{-}}\StringTok{ }\KeywordTok{setdiff}\NormalTok{(}\KeywordTok{seq\_len}\NormalTok{(n), train\_idx)}

\NormalTok{spectra\_tr \textless{}{-}}\StringTok{ }\NormalTok{spectra[train\_idx]}
\NormalTok{spectra\_te \textless{}{-}}\StringTok{ }\NormalTok{spectra[test\_idx]}
\KeywordTok{class}\NormalTok{(spectra\_tr) \textless{}{-}}\StringTok{ "spectra"}
\KeywordTok{class}\NormalTok{(spectra\_te) \textless{}{-}}\StringTok{ "spectra"}
\NormalTok{y\_tr \textless{}{-}}\StringTok{ }\NormalTok{group[train\_idx]}
\NormalTok{y\_te \textless{}{-}}\StringTok{ }\NormalTok{group[test\_idx]}
\end{Highlighting}
\end{Shaded}

\hypertarget{fit-a-single-model}{%
\section{Fit a single model}\label{fit-a-single-model}}

\texttt{fit\_mdm()} deconvolutes the spectra, aligns them, builds a
feature matrix and fits a lasso model via \texttt{cv.glmnet()}. The
model fitting grid (\texttt{mog}) controls the preprocessing parameters
explored. With a single-row grid, \texttt{fit\_mdm()} fits one model:

\begin{Shaded}
\begin{Highlighting}[]
\NormalTok{mog\_default \textless{}{-}}\StringTok{ }\KeywordTok{data.frame}\NormalTok{(}
    \DataTypeTok{nfit =} \DecValTok{3}\NormalTok{, }\DataTypeTok{smit =} \DecValTok{2}\NormalTok{, }\DataTypeTok{smws =} \DecValTok{5}\NormalTok{, }\DataTypeTok{delta =} \FloatTok{6.4}\NormalTok{, }\DataTypeTok{npmax =} \DecValTok{0}\NormalTok{,}
    \DataTypeTok{maxShift =} \DecValTok{100}\NormalTok{, }\DataTypeTok{maxCombine =} \DecValTok{30}
\NormalTok{)}
\NormalTok{mdm\_default \textless{}{-}}\StringTok{ }\KeywordTok{fit\_mdm}\NormalTok{(}
\NormalTok{    spectra\_tr, y\_tr, }\DataTypeTok{mog =}\NormalTok{ mog\_default, }\DataTypeTok{verbosity =} \DecValTok{0}
\NormalTok{)}
\KeywordTok{print}\NormalTok{(mdm\_default)}
\end{Highlighting}
\end{Shaded}

\begin{verbatim}
## metabodecon model (mdm)
##   npmax:         0
##   nfit:          3
##   smit:          2
##   smws:          5
##   delta:         6.4
##   maxShift:      100
##   maxCombine:    30
##   grid rows:     1
\end{verbatim}

\hypertarget{tune-preprocessing-via-grid-search}{%
\section{Tune preprocessing via grid
search}\label{tune-preprocessing-via-grid-search}}

A multi-row \texttt{mog} makes \texttt{fit\_mdm()} evaluate a grid of
preprocessing parameter combinations. For each row it builds the feature
matrix, runs \texttt{cv.glmnet()} with a fixed fold assignment, and
records the held-out accuracy and AUC at \texttt{lambda.min}. We use a
deliberately small grid here to keep vignette build time short:

\begin{Shaded}
\begin{Highlighting}[]
\NormalTok{mog \textless{}{-}}\StringTok{ }\KeywordTok{expand.grid}\NormalTok{(}
    \DataTypeTok{smit =} \DecValTok{2}\NormalTok{, }\DataTypeTok{smws =} \KeywordTok{c}\NormalTok{(}\DecValTok{3}\NormalTok{, }\DecValTok{5}\NormalTok{), }\DataTypeTok{delta =} \KeywordTok{c}\NormalTok{(}\DecValTok{4}\NormalTok{, }\DecValTok{6}\NormalTok{), }\DataTypeTok{nfit =} \DecValTok{5}\NormalTok{, }\DataTypeTok{npmax =} \DecValTok{0}\NormalTok{,}
    \DataTypeTok{maxShift =} \DecValTok{100}\NormalTok{, }\DataTypeTok{maxCombine =} \DecValTok{30}\NormalTok{,}
    \DataTypeTok{stringsAsFactors =} \OtherTok{FALSE}
\NormalTok{)}
\NormalTok{mdm\_tuned \textless{}{-}}\StringTok{ }\KeywordTok{fit\_mdm}\NormalTok{(}
\NormalTok{    spectra\_tr, y\_tr, }\DataTypeTok{mog =}\NormalTok{ mog, }\DataTypeTok{nfolds =} \DecValTok{5}\NormalTok{, }\DataTypeTok{verbosity =} \DecValTok{0}
\NormalTok{)}
\end{Highlighting}
\end{Shaded}

The performance grid shows how each parameter combination scored:

\begin{Shaded}
\begin{Highlighting}[]
\NormalTok{pg \textless{}{-}}\StringTok{ }\NormalTok{mdm\_tuned}\OperatorTok{$}\NormalTok{mog}
\NormalTok{pg \textless{}{-}}\StringTok{ }\NormalTok{pg[}\KeywordTok{order}\NormalTok{(}\OperatorTok{{-}}\NormalTok{pg}\OperatorTok{$}\NormalTok{auc), ]}
\NormalTok{knitr}\OperatorTok{::}\KeywordTok{kable}\NormalTok{(}
    \KeywordTok{head}\NormalTok{(pg, }\DecValTok{10}\NormalTok{),}
    \DataTypeTok{digits =} \KeywordTok{c}\NormalTok{(}\DecValTok{0}\NormalTok{, }\DecValTok{0}\NormalTok{, }\DecValTok{0}\NormalTok{, }\DecValTok{0}\NormalTok{, }\DecValTok{0}\NormalTok{, }\DecValTok{0}\NormalTok{, }\DecValTok{0}\NormalTok{, }\DecValTok{3}\NormalTok{, }\DecValTok{4}\NormalTok{),}
    \DataTypeTok{row.names =} \OtherTok{FALSE}\NormalTok{,}
    \DataTypeTok{caption =} \StringTok{"Top 10 parameter combinations by AUC."}
\NormalTok{)}
\end{Highlighting}
\end{Shaded}

\begin{longtable}[]{@{}rrrrrrrrr@{}}
\caption{Top 10 parameter combinations by AUC.}\tabularnewline
\toprule
smit & smws & delta & nfit & npmax & maxShift & maxCombine & acc &
auc\tabularnewline
\midrule
\endfirsthead
\toprule
smit & smws & delta & nfit & npmax & maxShift & maxCombine & acc &
auc\tabularnewline
\midrule
\endhead
2 & 5 & 4 & 5 & 0 & 100 & 30 & 0.776 & 0.8824\tabularnewline
2 & 3 & 6 & 5 & 0 & 100 & 30 & 0.701 & 0.7166\tabularnewline
2 & 5 & 6 & 5 & 0 & 100 & 30 & 0.672 & 0.6916\tabularnewline
2 & 3 & 4 & 5 & 0 & 100 & 30 & 0.567 & 0.5793\tabularnewline
\bottomrule
\end{longtable}

\hypertarget{compare-deconvolution-results}{%
\section{Compare deconvolution
results}\label{compare-deconvolution-results}}

Let us compare the deconvolution of the first training spectrum using
the default parameters and the best parameters found by the grid search:

\begin{Shaded}
\begin{Highlighting}[]
\CommentTok{\# Deconvolute first two training spectra with default and tuned parameters}
\NormalTok{pair \textless{}{-}}\StringTok{ }\NormalTok{spectra\_tr[}\DecValTok{1}\OperatorTok{:}\DecValTok{2}\NormalTok{]}
\KeywordTok{class}\NormalTok{(pair) \textless{}{-}}\StringTok{ "spectra"}

\NormalTok{d\_default \textless{}{-}}\StringTok{ }\KeywordTok{deconvolute}\NormalTok{(}
\NormalTok{    pair, }\DataTypeTok{nfit =} \DecValTok{3}\NormalTok{, }\DataTypeTok{smit =} \DecValTok{2}\NormalTok{, }\DataTypeTok{smws =} \DecValTok{5}\NormalTok{, }\DataTypeTok{delta =} \DecValTok{4}\NormalTok{, }\DataTypeTok{npmax =} \DecValTok{0}\NormalTok{, }\DataTypeTok{verbose =} \OtherTok{FALSE}
\NormalTok{)}

\NormalTok{m \textless{}{-}}\StringTok{ }\NormalTok{mdm\_tuned}\OperatorTok{$}\NormalTok{params}
\NormalTok{d\_tuned \textless{}{-}}\StringTok{ }\KeywordTok{deconvolute}\NormalTok{(}
\NormalTok{    pair, }\DataTypeTok{nfit =}\NormalTok{ m}\OperatorTok{$}\NormalTok{nfit, }\DataTypeTok{smit =}\NormalTok{ m}\OperatorTok{$}\NormalTok{smit, }\DataTypeTok{smws =}\NormalTok{ m}\OperatorTok{$}\NormalTok{smws,}
    \DataTypeTok{delta =}\NormalTok{ m}\OperatorTok{$}\NormalTok{delta, }\DataTypeTok{npmax =}\NormalTok{ m}\OperatorTok{$}\NormalTok{npmax, }\DataTypeTok{verbose =} \OtherTok{FALSE}
\NormalTok{)}

\KeywordTok{par}\NormalTok{(}\DataTypeTok{mfrow =} \KeywordTok{c}\NormalTok{(}\DecValTok{1}\NormalTok{, }\DecValTok{2}\NormalTok{), }\DataTypeTok{mar =} \KeywordTok{c}\NormalTok{(}\DecValTok{4}\NormalTok{, }\DecValTok{4}\NormalTok{, }\DecValTok{2}\NormalTok{, }\DecValTok{1}\NormalTok{))}
\KeywordTok{plot\_spectrum}\NormalTok{(d\_default[[}\DecValTok{1}\NormalTok{]], }\DataTypeTok{main =} \StringTok{"Default params"}\NormalTok{)}
\KeywordTok{plot\_spectrum}\NormalTok{(d\_tuned[[}\DecValTok{1}\NormalTok{]], }\DataTypeTok{main =} \StringTok{"Tuned params"}\NormalTok{)}
\end{Highlighting}
\end{Shaded}

\includegraphics[width=1\linewidth]{MDM_files/figure-latex/compare-decon-1}

\hypertarget{inspect-lasso-features}{%
\section{Inspect lasso features}\label{inspect-lasso-features}}

The non-zero coefficients of the lasso model reveal which chemical shift
positions are most important for distinguishing the two groups:

\begin{Shaded}
\begin{Highlighting}[]
\NormalTok{cf \textless{}{-}}\StringTok{ }\KeywordTok{coef}\NormalTok{(mdm\_tuned)}
\NormalTok{cf\_nz \textless{}{-}}\StringTok{ }\NormalTok{cf[cf[, }\DecValTok{1}\NormalTok{] }\OperatorTok{!=}\StringTok{ }\DecValTok{0}\NormalTok{, , drop =}\StringTok{ }\OtherTok{FALSE}\NormalTok{]}
\NormalTok{knitr}\OperatorTok{::}\KeywordTok{kable}\NormalTok{(}
    \KeywordTok{data.frame}\NormalTok{(}
        \DataTypeTok{feature =} \KeywordTok{rownames}\NormalTok{(cf\_nz),}
        \DataTypeTok{coefficient =} \KeywordTok{round}\NormalTok{(cf\_nz[, }\DecValTok{1}\NormalTok{], }\DecValTok{4}\NormalTok{)}
\NormalTok{    ),}
    \DataTypeTok{row.names =} \OtherTok{FALSE}\NormalTok{,}
    \DataTypeTok{caption =} \StringTok{"Non{-}zero lasso coefficients at lambda.min."}
\NormalTok{)}
\end{Highlighting}
\end{Shaded}

\begin{longtable}[]{@{}lr@{}}
\caption{Non-zero lasso coefficients at lambda.min.}\tabularnewline
\toprule
feature & coefficient\tabularnewline
\midrule
\endfirsthead
\toprule
feature & coefficient\tabularnewline
\midrule
\endhead
(Intercept) & 9.9722\tabularnewline
3.5492 & 1.1205\tabularnewline
3.52415 & 0.0571\tabularnewline
3.5231 & 0.0001\tabularnewline
3.5156 & 0.0002\tabularnewline
3.46535 & 0.0001\tabularnewline
3.4535 & 0.0000\tabularnewline
3.44075 & -0.0006\tabularnewline
3.43295 & -0.0155\tabularnewline
3.4196 & 0.0005\tabularnewline
3.3998 & 0.0015\tabularnewline
3.3902 & 0.0000\tabularnewline
3.38735 & 0.0000\tabularnewline
3.3818 & -0.0004\tabularnewline
\bottomrule
\end{longtable}

\hypertarget{predict-test-samples}{%
\section{Predict test samples}\label{predict-test-samples}}

We use \texttt{predict.mdm()} to classify the held-out test spectra.
This internally deconvolutes and aligns them against the reference
spectrum stored in the model:

\begin{Shaded}
\begin{Highlighting}[]
\NormalTok{preds \textless{}{-}}\StringTok{ }\KeywordTok{predict}\NormalTok{(mdm\_tuned, spectra\_te, }\DataTypeTok{type =} \StringTok{"all"}\NormalTok{, }\DataTypeTok{verbosity =} \DecValTok{0}\NormalTok{)}
\NormalTok{results \textless{}{-}}\StringTok{ }\KeywordTok{data.frame}\NormalTok{(}
    \DataTypeTok{sample =} \KeywordTok{get\_names}\NormalTok{(spectra\_te),}
    \DataTypeTok{true =}\NormalTok{ y\_te,}
    \DataTypeTok{predicted =}\NormalTok{ preds}\OperatorTok{$}\NormalTok{class,}
    \DataTypeTok{prob\_B =} \KeywordTok{round}\NormalTok{(preds}\OperatorTok{$}\NormalTok{prob, }\DecValTok{3}\NormalTok{)}
\NormalTok{)}
\end{Highlighting}
\end{Shaded}

\begin{Shaded}
\begin{Highlighting}[]
\NormalTok{acc \textless{}{-}}\StringTok{ }\KeywordTok{mean}\NormalTok{(preds}\OperatorTok{$}\NormalTok{class }\OperatorTok{==}\StringTok{ }\NormalTok{y\_te)}
\NormalTok{auc \textless{}{-}}\StringTok{ }\NormalTok{metabodecon}\OperatorTok{:::}\KeywordTok{AUC}\NormalTok{(y\_te, preds}\OperatorTok{$}\NormalTok{prob)}
\KeywordTok{cat}\NormalTok{(}\KeywordTok{sprintf}\NormalTok{(}\StringTok{"Test accuracy: \%.1f\%\%}\CharTok{\textbackslash{}n}\StringTok{"}\NormalTok{, acc }\OperatorTok{*}\StringTok{ }\DecValTok{100}\NormalTok{))}
\end{Highlighting}
\end{Shaded}

\begin{verbatim}
## Test accuracy: 72.7%
\end{verbatim}

\begin{Shaded}
\begin{Highlighting}[]
\KeywordTok{cat}\NormalTok{(}\KeywordTok{sprintf}\NormalTok{(}\StringTok{"Test AUC:      \%.4f}\CharTok{\textbackslash{}n}\StringTok{"}\NormalTok{, auc))}
\end{Highlighting}
\end{Shaded}

\begin{verbatim}
## Test AUC:      0.8824
\end{verbatim}

\begin{Shaded}
\begin{Highlighting}[]
\NormalTok{knitr}\OperatorTok{::}\KeywordTok{kable}\NormalTok{(}
    \KeywordTok{table}\NormalTok{(}\DataTypeTok{True =}\NormalTok{ y\_te, }\DataTypeTok{Predicted =}\NormalTok{ preds}\OperatorTok{$}\NormalTok{class),}
    \DataTypeTok{caption =} \StringTok{"Confusion matrix on test set."}
\NormalTok{)}
\end{Highlighting}
\end{Shaded}

\begin{longtable}[]{@{}lrr@{}}
\caption{Confusion matrix on test set.}\tabularnewline
\toprule
& A & B\tabularnewline
\midrule
\endfirsthead
\toprule
& A & B\tabularnewline
\midrule
\endhead
A & 15 & 1\tabularnewline
B & 8 & 9\tabularnewline
\bottomrule
\end{longtable}

\hypertarget{benchmark-with-nested-cross-validation}{%
\section{Benchmark with nested
cross-validation}\label{benchmark-with-nested-cross-validation}}

To obtain an unbiased estimate of end-to-end predictive performance,
\texttt{benchmark()} wraps \texttt{fit\_mdm()} in an outer
cross-validation loop. The outer held-out fold is never seen during
preprocessing or model fitting, so the resulting accuracy and AUC are
honest estimates.

A typical call looks like this:

\begin{Shaded}
\begin{Highlighting}[]
\NormalTok{bm \textless{}{-}}\StringTok{ }\KeywordTok{benchmark}\NormalTok{(}
\NormalTok{    spectra\_tr, y\_tr,}
    \DataTypeTok{fun =} \StringTok{"fit\_mdm"}\NormalTok{, }\DataTypeTok{k =} \DecValTok{3}\NormalTok{, }\DataTypeTok{nworkers =} \DecValTok{4}\NormalTok{,}
    \DataTypeTok{mog =}\NormalTok{ mog, }\DataTypeTok{nfolds =} \DecValTok{5}
\NormalTok{)}
\CommentTok{\# bm$predictions contains per{-}sample held{-}out predictions}
\CommentTok{\# bm$models contains the fit\_mdm model for each outer fold}
\KeywordTok{cat}\NormalTok{(}\KeywordTok{sprintf}\NormalTok{(}
    \StringTok{"Nested CV accuracy: \%.1f\%\%}\CharTok{\textbackslash{}n}\StringTok{"}\NormalTok{,}
    \KeywordTok{mean}\NormalTok{(bm}\OperatorTok{$}\NormalTok{predictions}\OperatorTok{$}\NormalTok{true }\OperatorTok{==}\StringTok{ }\NormalTok{bm}\OperatorTok{$}\NormalTok{predictions}\OperatorTok{$}\NormalTok{pred) }\OperatorTok{*}\StringTok{ }\DecValTok{100}
\NormalTok{))}
\end{Highlighting}
\end{Shaded}

Note that \texttt{benchmark()} is computationally expensive because it
runs the full grid search for each outer fold. It is therefore not
executed in this vignette.

\end{document}
